\section{Method}
\subsection{Problem Formulation}
In offline RL, the agent is provided with a fixed dataset of previously collected trajectories.
%Let $\mathcal{D}=\{\boldsymbol{\tau}^i\}_{i=1}^N$ be an dataset of trajectories. Each trajectory is a finite sequence of transitions and terminal marker, i.e. it has the form of (\ref{eq:trace_form}) plus a special symbol $\texttt{END}$ at the end of the trace. At each timestamp, together with state dimensions, action dimensions and rewards, we can have an optional environment-specific auxiliary fields (e.g., safety cost).
Let $\mathcal{D}=\{\boldsymbol{\tau}^i\}_{i=1}^N$ be an dataset of trajectories. Each trajectory is a finite sequence of transitions and an \emph{end-of-trace} symbol $\texttt{EOT}$, i.e. $\boldsymbol{\tau} = (y^1,\ldots,y^T,\texttt{EOT})$, where each transition token $y^j$ group encodes state, action, reward, and optional environment-specific auxiliary fields (e.g., safety cost). 
In addition to the return signal, here we add also some prior knowledge about the task, expressed as \ltlf constraints, to capture high-level behavioral requirements. Given an \ltlf formula $\varphi$, our objective is to learn a trajectory generation policy that maximizes both the expected return of the policy and the probability that generated trajectories satisfy $\varphi$.
%
We define atomic propositions through an environment-specific extractor 
\[
\Pi(y^t)\in\{0,1\}^K
\]
%
\noindent
and the \ltlf formula $\varphi$ is defined over these propositions. Note that, in our experiments, proposition extraction is primarily \textit{state-centric}, meaning that $\varphi$ is defined only over state dimensions, but our approach is fully general and supports also imposing constraints on action, state-action pairs and auxiliary-terms.

A logic adapter is used to define the authoritative token schema and the unique end token identifier shared by dataset builders, token-to-symbol mapping, DFA rollout, and logic-loss evaluation. Full episodes contain exactly one explicit end row; sliding training windows may omit the end marker if they are strict subsegments. Canonical satisfaction evaluation is defined on complete finite traces, and traces without an end marker are treated as incomplete and unsatisfied in canonical mode. This design avoids inconsistent termination handling across modules and ensures agreement between hard DFA evaluation and differentiable soft evaluation.

As explained in the background, we assume an autoregressive neural model $f_\theta$ with trainable parameters $\theta$ that estimates an approximation $P_\theta$ of the probability of the next event $\tau_t$ given a trace of previous events $\boldsymbol{\tau}_{<t}$ (Equation~\ref{eq:next_act_pred}):
\begin{equation} 
\label{eq:RNN}
%\[
\begin{array}{ll}
     \tilde{y}_t = f_\theta(\boldsymbol{\tau}_{<t}) \\
     P(\tau_{t} = \tau_i \mid \boldsymbol{\tau}_{<t}) \approx \tilde{y}_t[i].
\end{array}
%\]
\end{equation}
Note that we do not make any assumptions about the neural model, except that it can estimate the probability of the next activity given a sequence of previous ones. As a result, our approach is entirely \textit{model-agnostic} and can be readily applied to any autoregressive model (providing the fact that the logic adapter is changed consistently). The model parameters are typically trained using a supervised loss $L_{\mathcal{D}}$, evaluated on a dataset $\mathcal{D}$ of ground-truth traces obtained by observing the process. The loss for a trace $\boldsymbol{\tau} \in \mathcal{D}$ of length $T$ is defined as follows:
\begin{equation} \label{eq:sup_loss}
    L_{\mathcal{D}}(\boldsymbol{\tau}) = \frac{1}{T} \sum_{t=1}^T \text{cross-entropy}(f_\theta(\boldsymbol{\tau}_{<t}), \tau_{t}) .
\end{equation}
This loss trains the network to predict the next symbol in a trace so as to closely mimic the data in the dataset.

Here, the goal is 
%Our goal is for the language generated by the autoregressor $f_\theta$ to be strictly contained within the language of strings accepted by the \ltlf formula $\varphi$, denoted as $\calL(A_{\varphi})$. However, the language produced by the network is \emph{unbounded}, as it is only \emph{softly assigned}: in other words, the network can generate any possible string, each with a different probability.
%
%Our method therefore aims 
to maximize the probability $P_{\theta \vDash \varphi}$ that traces $\boldsymbol{\tau} \sim P_\theta$, sampled from the autoregressor, satisfy the specification:
\begin{equation} \label{eq:logic_prob}
P_{\theta \vDash \varphi} = \mathbb{E}_{\boldsymbol{\tau} \sim P_\theta}[\boldsymbol{\tau} \vDash \varphi] = 
\sum_{\boldsymbol{\tau}} P_\theta(\boldsymbol{\tau}) \, \mathbb{I}\{\boldsymbol{\tau} \vDash \varphi\}.
\end{equation}

\begin{comment}
Note that, in order to compute the exact probability of knowledge satisfaction, one would need to enumerate \textit{all possible suffixes} of maximum length $T$ and sum their acceptance probabilities. However, this set has exponential size, making exact computation unfeasible. Furthermore, we aim to impose the maximization of this probability as a \textit{training objective}. Therefore, our goal is to design a fast and differentiable procedure to approximate it at each optimization step of the autoregressor.

We propose two logic loss functions: a local, activity-level loss $L_\varphi^{\text{loc}}$, and a global, trace-level loss $L_\varphi^{\text{glob}}$. We show that both contribute positively to the learning process, demonstrating the effectiveness of integrating \ltlf knowledge into autoregressive training. The choice between the two logic losses depends on the type of \ltlf knowledge available. The global loss $L_\varphi^{\text{glob}}$ can always be applied, regardless of the structure of the formula $\varphi$, but it is more computationally demanding, as it requires generating entire suffixes during training.

In general, the compliance of a trace with $\varphi$ can only be assessed on complete traces. Indeed, whether the current partial prediction $\boldsymbol{a}^{\leq t}$ satisfies (or violates) $\varphi$ does not guarantee that the final predicted trace $\boldsymbol{a}$ will also satisfy (or violate) it. This introduces a challenge in evaluating \ltlf constraints, as opposed to approaches that rely on local constraints~\cite{montecarlo_llm}, which can be verified step-by-step during generation.

However, some types of formulas, such as safety constraints, once translated into a DFA, contain \emph{failing states}, which can be used to guide generation locally. A failing state is a state in the DFA from which no accepting state is reachable. When a partial trace enters such a state, it has irreversibly violated the formula, and no continuation can satisfy the knowledge. Our local loss $L_\varphi^{\text{loc}}$ exploits the presence of failing states to guide the autoregressor at every generation step. As a result, it provides richer feedback at the activity level, but it can only be used with formulas that admit such a DFA structure. \textcolor{blue}{Nevertheless, this limitation is overcome during the DFA preprocessing phase, where a failing state is always added to handle the \texttt{EOT} symbol, as discussed in Section~\ref{sec:knowledge_preprocessing}.}
\end{comment}

\noindent
Computing this probability exactly is infeasible, since it would require to enumerate al possible traces of maximum length $T$. In \cite{UmiliPMAI24,mezini2025neuro} a differentiable procedure to approximate the previous probability at each optimization step of the autoregressor is designed. In particular, they introduce a logic loss function $L_\varphi$, and compute the training objective as a linear combination of the logic and the supervised loss:
\begin{equation}
\label{eq:loss_balance}
    L = \alpha L_{\mathcal{D}} + (1 - \alpha) L_\varphi,
\end{equation}
%
with $\alpha$ being a trade-off parameter between 0 and 1 that balances the influence of each loss on the training process.
Here we adapt the logic loss computation to offline RL domains and apply it at training time, together with other techniques at test time to maximize adherence with the constraints. 


\subsection{Automata-Based Satisfaction and Logic Loss}

Now, we can briefly discuss how it is possible to approximate the target probability in Equation \ref{eq:logic_prob}. Specifically, it is possible to use a Monte Carlo estimation by sampling a set of complete traces $\{\boldsymbol{\tau}^{(1)}, \boldsymbol{\tau}^{(2)}, \ldots, \boldsymbol{\tau}^{(N)}\} \sim P_\theta$ according to the distribution learned by the autoregressor, and compute an approximation of the target probability as follows:
\begin{equation} \label{eq:montecarlo_approx}
    \hat{P}_{\theta \vDash \varphi} = \frac{1}{N} \sum_{i=1}^N \mathbb{I} \{ \boldsymbol{\tau}^{(i)} \vDash \varphi \}.
\end{equation}
%
Note that the sampled traces $\boldsymbol{\tau}^{(i)}$ are generated by a NN that returns a probability distribution at each time step, and thus they are replaced by a probabilistic counterpart $\tilde{\boldsymbol{\tau}}^{(i)}$, where symbols are sampled from that probability distribution. Moreover, the indicator function $\mathbb{I}$ in Equation~\ref{eq:montecarlo_approx} is replaced by a method that computes the (probabilistic) compliance of a probabilistic trace with the knowledge. To achieve this, it is possible to leverage on \emph{(i)} the Gumbel-Softmax sampling, to generate differentiable, near one-hot suffixes during training and \emph{(ii)} DeepDFA \cite{deepdfa_ecai2024}, a neurosymbolic framework that encodes temporal logic properties as a recurrent layer, enabling efficient and differentiable evaluation of logical constraints.

\paragraph{Gumbel-Softmax sampling.} In order to sample from the probability distribution returned by a NN, but still maintain differentiability, we can use the Gumbel-Softmax reparametrization trick, which produces approximations of discrete symbols as one-hot-like vectors. Given the predicted token distribution $\tilde{y}_t$, we get 
\begin{equation}
\tilde{x}_t = \text{softmax} \left( \frac{\log (\tilde{y}_t) +G}{temp} \right)
\end{equation}
%
\noindent
where $G$ is a Gumbel noise vector and $temp$ is a temperature parameter controlling the sharpness of the distribution. These soft tokens are then fed into the differentiable DFA, which computes the probability that the sampled trajectory satisfies the LTLf specification. As $temp \rightarrow 0$, the output approaches a discrete one-hot vector, while for $temp = 1$, it remains close to the original continuous probabilities in $\tilde{y}_t$. Since the next activity is only \emph{probabilistically} grounded, we denote it as $\tilde{\tau}_t$.

\paragraph{DeepDFA.}
Each \ltlf constraint $\varphi$ can be compiled into an equivalent DFA 
\[
A_\varphi = (\Sigma, Q, q_0, \delta, F)
\]
%
\noindent
where $\Sigma$ is the alphabet, $Q$ is the set of states, $q_0$ is the initial state, $\delta$ is the transition function, and $F$ is the set of accepting states. There exists automatic translation tools, such as \texttt{ltlf2DFA}~\cite{fuggitti-ltlf2dfa}, that can be used off-the-shelf. We extend the DFA to handle the special $\texttt{EOT}$ symbol, such that only traces including $\texttt{EOT}$ are accepted. Given a trace of symbols $\boldsymbol{\sigma}$, the automaton processes the corresponding sequence of symbols and determines whether the trajectory satisfies the specification:
\[
\boldsymbol{\sigma} \models \varphi \text{ if and only if } \delta^*(q_0, \boldsymbol{\sigma}) \in F
\]

\textcolor{blue}{***}

We expose two aligned satisfaction interfaces:

\paragraph{Hard path.}
Given token IDs, we map each timestep to a symbol ID and roll the DFA exactly, returning Boolean acceptance per trajectory.

\paragraph{Soft path.}
Given token probability distributions, we aggregate token probabilities into symbol probabilities and perform a soft DFA rollout to obtain an acceptance probability in $[0,1]$:
\[
P_t(s)=\sum_{\textcolor{blue}{k:m_t(k)=s}} p_t(k), \qquad \textcolor{blue}{\hat{S}_\varphi(\tau)\in[0,1]}.
\]
\textcolor{blue}{Ele: define things in blue in previous equation}
One-hot token distributions recover the hard-path result up to numerical tolerance.

\textcolor{blue}{***}

Many alternatives exist in the literature to evaluate continuous satisfaction of temporal constraints over sequences of \textit{soft} (probabilistic or fuzzy) symbol assignments \cite{umili_kr23,deepdfa_ecai2024,DonadelloFIMM25,NesyA}. In our work we used DeepDFA, following what has been done in \cite{UmiliPMAI24,mezini2025neuro}. DeepDFA is a neural, probabilistic relaxation of a standard DFA, where the automaton is represented in matrix form and the input symbols, states, and outputs are \textit{probabilistically grounded}. This allows us to compute the probabilistic compliance $P_{\mathrm{DDFA}}(\tilde{\boldsymbol{\tau}}^{(i)} \vDash \varphi)$ of a sampled trace $\tilde{\boldsymbol{\tau}}^{(i)}$ with the knowledge $\varphi$. Therefore, Equation~\ref{eq:montecarlo_approx} becomes:
    \begin{equation} \label{eq:montecarlo_approx_ddfa}
    \hat{P}_{\theta \vDash \varphi} = \frac{1}{N} \sum_{i=1}^N P_{\mathrm{DDFA}}(\tilde{\boldsymbol{\tau}}^{(i)} \vDash \varphi).
\end{equation}
%
\noindent
Finally, we are able to compute the logic loss $L_\varphi$, which enforces the satisfaction of prior knowledge over entire traces:
\begin{equation}
L_\varphi = - \log \left( \hat{P}_{\theta \vDash \varphi} \right).
\end{equation}
and we use it in Equation \ref{eq:loss_balance}.

\begin{comment}
\subsection{TT: Logic-Regularized Training and Automaton-Aware Decoding}
TT models trajectories autoregressively. We use standard likelihood training and optionally add a logic regularization term weighted by $\alpha$ (swept in experiments), using soft satisfaction computed from sampled/generated traces. We also support multiple DFA composition modes (\texttt{single}, \texttt{product}, \texttt{multi}) for composite specifications.

At evaluation, we compare:
\begin{itemize}
\item greedy decoding,
\item unconstrained beam search,
\item automaton-aware constrained beam search.
\end{itemize}
Each beam maintains a DFA state. We support soft satisfaction reranking and hard pruning of beams that enter identified rejecting sink states.

\subsection{DT: Baseline and Logic-Enhanced Variants}
We consider the following DT family methods:
\begin{itemize}
\item \textbf{DT vanilla:} greedy inference baseline.
\item \textbf{DT-LTF-1:} training-time logic regularization with rollout-based soft satisfaction (logic weight \texttt{logic\_alpha} and rollout controls).
\item \textbf{DT-LTF-2:} constrained test-time action selection (\texttt{dt\_mode constrained}) using candidate actions, short-horizon lookahead, DFA state updates, and optional hard pruning.
\item \textbf{DT-LTF-3:} offline kNN suffix continuation planner (\texttt{dt\_mode knn}) scoring retrieved continuations with return and satisfaction proxies.
\end{itemize}
For DT evaluation, we additionally report a random policy baseline as a reference.
\end{comment}
